\documentclass[12pt]{article}

\usepackage[utf8]{inputenc}
\usepackage{geometry}
\usepackage{booktabs}
\geometry{margin=1in}

\title{Communication and Coordination Plan}
\author{}
\date{}

\begin{document}

\maketitle

\section{Communication Channels}

Effective communication is essential to ensure coordination across all components of the system development process. The team will use the following channels:

\begin{itemize}
    \item \textbf{WhatsApp}: Primary channel for real-time communication, quick clarifications, and urgent coordination.
    
    \item \textbf{GitHub}: Main platform for version control, issue tracking, pull requests, and technical discussions. All development-related communication must be documented here.
    
    \item \textbf{Weekly Meetings}: Formal synchronization space to review progress, blockers, and next steps.
\end{itemize}

\section{Meeting Frequency and Structure}

\begin{itemize}
    \item \textbf{Weekly Status Meeting}:
    \begin{itemize}
        \item Frequency: Once per week
        \item Duration: 45–60 minutes
        \item Objective: Review progress, validate deliverables, identify risks, and plan next iteration
    \end{itemize}
    
    \item \textbf{Ad-hoc Meetings}:
    \begin{itemize}
        \item Triggered when critical issues arise or architectural decisions are required
        \item Short and focused (15–30 minutes)
    \end{itemize}
\end{itemize}

\section{Issue Escalation Process}

To ensure timely resolution of blockers, the following escalation protocol is defined:

\begin{enumerate}
    \item Issues are first reported and documented in \textbf{GitHub Issues}.
    \item If unresolved within 24 hours, the issue is escalated via \textbf{WhatsApp}.
    \item If the issue impacts deadlines or system stability, it is escalated to the \textbf{Project Lead}.
    \item Critical issues are addressed in an \textbf{ad-hoc meeting} with relevant team members.
\end{enumerate}

\section{Control and Monitoring Tools}

The team will use structured tools to track progress and ensure transparency:

\begin{itemize}
    \item \textbf{GitHub Projects}: Used for task management, sprint planning, and workflow visualization (Kanban board).
    
    \item \textbf{Burndown Charts}: Used to monitor sprint progress and detect deviations from planned timelines.
    
    \item \textbf{Pull Requests and Code Reviews}: Ensure code quality, traceability, and collaborative validation before merging changes.
    
    \item \textbf{Issue Tracking}: All bugs, features, and improvements must be registered and tracked through GitHub Issues.
\end{itemize}

\section{Summary}

This communication and coordination structure ensures alignment, transparency, and efficient problem resolution. By combining synchronous (meetings, WhatsApp) and asynchronous (GitHub) communication channels, the team can maintain both agility and traceability throughout the project lifecycle.

\end{document}